\begin{figure}[t]
\centering
\begin{tikzpicture}
\begin{axis}[edge, height=0.52\linewidth,
    xmin=1990, xmax=2022, xtick={1990,2000,2010,2022},
    x tick label style={/pgf/number format/1000 sep={}},
    ymin=0, ymax=36, ylabel={Adults with diabetes, \% (age-standardised)},
    legend style={at={(0.02,0.98)},anchor=north west}]
  \addplot[lineEmber] table[x=year,y=pacific] {\datadir/diabetes_groups.csv};
  \addlegendentry{Pacific island states (median)}
  \addplot[lineTeal] table[x=year,y=sids] {\datadir/diabetes_groups.csv};
  \addlegendentry{All small island states (median)}
  \addplot[lineSlate] table[x=year,y=all] {\datadir/diabetes_groups.csv};
  \addlegendentry{All countries (median)}
\end{axis}
\end{tikzpicture}
\caption{\takeaway{Island diabetes rates were high and kept rising.}
  The median Pacific state went from \diabPacificNinety\% to \diabPacificNow\%
  between 1990 and 2022; the median country from \diabAllNinety\% to
  \diabAllNow\%. The rest of the world is catching up, so the relative gap
  narrowed (from \diabPacificRatioNinety{} to \diabPacificRatioNow{} times) even
  as the absolute gap held. Mean of male and female estimates, adults 18+.
  \source{NCD-RisC \parencite{ncdrisc2024}; UN M49.}}
\label{fig:diabetes}
\end{figure}
